\documentclass[12pt]{article}
\usepackage[english, bulgarian]{babel}
\usepackage[utf8]{inputenc}
\usepackage[T2A]{fontenc}
\usepackage{amssymb}
\usepackage{hyperref, fancyhdr, lastpage, fancyvrb, tcolorbox, titlesec}
\usepackage{array, tabularx, colortbl}
\usepackage{tikz}
\usepackage{venndiagram}
\usepackage{amsthm, bm}
\usepackage{relsize}
\usepackage{amsmath,physics}
\usepackage{mathtools}
\usepackage{subcaption}
\usepackage{scalerel}
\usepackage{theoremref}
\usepackage{circuitikz}
\usepackage{geometry}
\usepackage{stmaryrd}
\usepackage{forest}
\usepackage{cancel}
\usepackage{faktor}
\usepackage{gensymb}
\usepackage[shortlabels]{enumitem}
\usepackage[symbol]{footmisc}
\usetikzlibrary{automata, arrows, positioning, shapes}
\useforestlibrary{linguistics}

\ExplSyntaxOn
\NewDocumentCommand{\opair}{m}
 {
  \langle\mspace{2mu}
  \clist_set:Nn \l_tmpa_clist { #1 }
  \clist_use:Nn \l_tmpa_clist {,\mspace{3mu plus 1mu minus 1mu}\allowbreak}
  \mspace{2mu}\rangle
}
\ExplSyntaxOff

\hypersetup{
    colorlinks=true,
    linktoc=all,
    linkcolor=blue
}

\setlength\parindent{0pt}

\newcommand{\N}{\mathbb{N}}
\newcommand{\Z}{\mathbb{Z}}
\newcommand{\Q}{\mathbb{Q}}
\newcommand{\R}{\mathbb{R}}
\newcommand{\E}{\mathbb{E}}

\newcommand{\vars}{\operatorname{Vars}}
\newcommand{\free}{\operatorname{Free}}

\newcommand{\logand}{\; \& \;}

\newcommand{\calA}{\mathcal{A}}
\newcommand{\calS}{\mathcal{S}}
\newcommand{\calD}{\mathcal{D}}
\newcommand{\calL}{\mathcal{L}}
\newcommand{\calZ}{\mathcal{Z}}
\newcommand{\calQ}{\mathcal{Q}}
\newcommand{\calR}{\mathcal{R}}
\newcommand{\calN}{\mathcal{N}}
\newcommand{\calM}{\mathcal{M}}
\newcommand{\calF}{\mathcal{F}}
\newcommand{\calP}{\mathcal{P}}
\newcommand{\calC}{\mathcal{C}}
\newcommand{\calG}{\mathcal{G}}
\newcommand{\calB}{\mathcal{B}}

\newcommand{\dequiv}{\stackrel{\text{деф.}}{\longleftrightarrow}}

\newcommand{\db}[1]{\llbracket #1 \rrbracket}

\DeclareMathOperator*{\bigand}{\scalerel*{\&}{\prod}}
\DeclareMathOperator*{\bigor}{\scalerel*{\lor}{\sum}}

\newtheorem*{definition}{Дефиниция}
\newtheorem*{claim}{Твърдение}
\newtheorem*{property}{Свойство}
\newtheorem*{hint}{Упътване}
\theoremstyle{definition}
\newtheorem{problem}{Задача}
\newtheorem*{solution}{Решение}
\newtheorem*{notation}{Нотация}

\begin{document}

\title{Кратки решения на задачите от семестриалното контролно по ЕАИ на специалност ``Компютърни науки'', проведено на 26 април 2026 г.}
\date{}

\maketitle

\begin{problem}[{\bf 5 точки}]
Постройте минимален ДКА за езика:
\[
  L = \{ a, b \}^\star \cdot \{ abb, ab \}.
\]
\end{problem}

\begin{solution}
Строим ДКА по метода на Бжозовски, като започваме с началното състояние:
\begin{align*}
  L_0 := L & = \calL \db{(a + b)^\star (abb + ab)} = \calL \db{((a + b)(a + b)^\star + \varepsilon) (abb + ab)} \\ 
           & = \calL \db{(a + b)(a + b)^\star (abb + ab) + abb + ab} \\
           & = \calL \db{a (a + b)^\star (abb + ab) + abb + ab + b (a + b)^\star (abb + ab)}.
\end{align*}
Пресмятаме преходите на $L_0$:
\begin{align*}
  a^{-1}(L_0) & = \calL \db{(a + b)^\star (abb + ab) + bb + b} := L_1, \text{ и } \\
  b^{-1}(L_0) & = \calL \db{(a + b)^\star (abb + ab)} = L_0.
\end{align*}
Тъй като $b \in L_1$ и $b \not\in L_0$, получаваме $L_1 \neq L_0$.

Сега пресмятаме преходите на $L_1$:
\begin{align*}
  a^{-1}(L_1) & = a^{-1}(\calL \db{(a + b)^\star (abb + ab) + bb + b}) \\ 
              & = a^{-1}(L_0) \cup a^{-1}(\calL \db{bb + b}) = L_1 \cup \varnothing = L_1, \text{ и } \\
  b^{-1}(L_1) & = b^{-1}(\calL \db{(a + b)^\star (abb + ab) + bb + b}) \\ 
              & = b^{-1}(L_0) \cup b^{-1}(\calL \db{bb + b}) \\ 
              & = \calL \db{(a + b)^\star (abb + ab) + b + \varepsilon} := L_2.
\end{align*}
Тъй като $\varepsilon \in L_2$ и $\varepsilon \not\in L_0, L_1$, имаме $L_2 \neq L_1, L_0$.

След това пресмятаме преходите на $L_2$:
\begin{align*}
  a^{-1}(L_2) & = a^{-1}(\calL \db{(a + b)^\star (abb + ab) + b + \varepsilon}) \\ 
              & = a^{-1}(L_0) \cup a^{-1}(\calL \db{b + \varepsilon}) = L_1 \cup \varnothing = L_1, \text{ и } \\
  b^{-1}(L_2) & = b^{-1}(\calL \db{(a + b)^\star (abb + ab) + b + \varepsilon}) \\ 
              & = b^{-1}(L_0) \cup b^{-1}(\calL \db{b + \varepsilon}) \\ 
              & = \calL \db{(a + b)^\star (abb + ab) + \varepsilon} := L_3.
\end{align*}
Понеже $\varepsilon \in L_3$ и $\varepsilon \not\in L_0, L_1$, имаме $L_3 \neq L_0, L_1$.
Освен това, тъй като $b \in L_2$ и $b \not\in L_3$, получаваме $L_3 \neq L_2$.

Остава да сметнем преходите на $L_3$:
\begin{align*}
  a^{-1}(L_3) & = a^{-1}(\calL \db{(a + b)^\star (abb + ab) + \varepsilon}) \\ 
              & = a^{-1}(L_0) \cup a^{-1}(\calL \db{\varepsilon}) = L_1 \cup \varnothing = L_1, \text{ и } \\
  b^{-1}(L_3) & = b^{-1}(\calL \db{(a + b)^\star (abb + ab) + \varepsilon}) \\ 
              & = b^{-1}(L_0) \cup b^{-1}(\calL \db{\varepsilon}) \\ 
              & = \calL \db{(a + b)^\star (abb + ab) + \varnothing} := L_0.
\end{align*}
Получаваме следния автомат:
\begin{center}
    \begin{tikzpicture}[shorten >=1pt,node distance=2.5cm,>=stealth',thick]
        \node[initial, state, initial text=] (1) {$L_0$};
        \node[state] [right of=1] (2) {$L_1$};
        \node[accepting, state] [below of=2] (3) {$L_2$};
        \node[accepting, state] [below of=1] (4) {$L_3$};
        \path[->] (1) edge [loop above] node[above] {$b$} (1);
        \path[->] (1) edge node[above] {$a$} (2);
        \path[->] (2) edge [loop above] node[above] {$a$} (2);
        \path[->] (2) edge [bend left] node[right] {$b$} (3);
        \path[->] (3) edge node[left] {$a$} (2);
        \path[->] (3) edge node[above] {$b$} (4);
        \path[->] (4) edge node[above] {$a$} (2);
        \path[->] (4) edge node[left] {$b$} (1);
    \end{tikzpicture}
\end{center}
\end{solution}

\newpage

\begin{problem}
За две думи $\alpha$ и $\beta$ с равна дължина над азбуката $\Sigma = \{ 0, 1 \}$, дефинираме операцията $\texttt{diff}(\alpha,\beta)$ по следния начин:
\begin{align*}
  & \texttt{diff}(\varepsilon, \varepsilon) = 0\\
  & \texttt{diff}(x \alpha, y \beta) = (x-y) + \texttt{diff}(\alpha, \beta).
\end{align*}
Например, $\texttt{diff}(010, 101) = -1$.
Сега дефинираме операцията върху езици над азбуката $\Sigma = \{ 0, 1 \}$ по следния начин:
\[
  \texttt{diff}(L) = \{ \alpha \in \Sigma^\star \mid (\exists \beta \in L)[|\alpha| = |\beta| \logand \texttt{diff}(\alpha,\beta) = 0] \}.
\]

\begin{enumerate}[(a)]
  \item Дайте пример за безкраен регулярен език $L$, за който езикът $\texttt{diff}(L)$ е регулярен. {\bf (2 точки)}
  \item Дайте пример за регулярен език $L$, за който езикът $\texttt{diff}(L)$ не е регулярен. {\bf (3 точки)}
\end{enumerate}
\end{problem}

\begin{solution}
Първо можем да забележим, че за всяко $\alpha, \beta \in \Sigma^\star$ с $|\alpha| = |\beta| = n$:
\begin{align*}
  \texttt{diff}(\alpha, \beta) & = \sum\limits_{i = 1}^n (\alpha[i] - \beta[i]) \\ 
                               & = \sum\limits_{i = 1}^n \alpha[i] - \sum\limits_{i = 1}^n \beta[i] = |\alpha|_1 - |\beta|_1.
\end{align*}
Тогава, ако $\texttt{diff}(\alpha, \beta) = 0$, то $|\alpha|_1 = |\beta|_1$. Тъй като $|\alpha| = |\beta|$, имаме и $|\alpha|_0 = |\beta|_0$.
В частност:
\begin{itemize}
  \item ако $\texttt{diff}(\alpha, 0^n) = 0$, то $\alpha = 0^n$;
  \item ако $\texttt{diff}(\alpha, (01)^n) = 0$, то $|\alpha|_0 = |\alpha|_1 = n$.
\end{itemize}
Така:
\begin{enumerate}[(a)]
  \item $\texttt{diff}(\underbrace{\{ 0 \}^\star}_{\text{безкраен регулярен}}) = \underbrace{\{ 0 \}^\star}_{\text{регулярен}}$;
  \item $\texttt{diff}(\underbrace{\{ 01 \}^\star}_{\text{регулярен}}) = \underbrace{\{ \alpha \in \Sigma^\star \mid |\alpha|_0 = |\alpha|_1 \}}_{\text{класически нерегулярен}}$ \footnote[2]{Ако пресечем $\{ \alpha \in \Sigma^\star \mid |\alpha|_0 = |\alpha|_1 \}$ с регулярния $\{ 0 \}^\star \cdot \{ 1 \}^\star$, ще получим $\{ 0^n 1^n \mid n \in \N \}$, който е нерегулярен. Затова $\{ \alpha \in \Sigma^\star \mid |\alpha|_0 = |\alpha|_1 \}$ не е регулярен.}.
\end{enumerate}
\end{solution}

\newpage

\begin{problem}[{\bf 5 точки}]
За произволен език $L$ над азбуката $\Sigma = \{ a, b \}$ дефинираме операцията:
\[
  \texttt{Part}(L) = \{ \beta \in \Sigma^\star \mid (\exists \alpha \in \Sigma^\star)[\alpha \beta \cdot \alpha \beta \in L] \}.
\]
Докажете, че ако $L$ е регулярен език, то $\texttt{Part}(L)$ е регулярен език.
\end{problem}

\begin{solution}
Този език много напомня на разглеждания на лекции/упражнения език (полагайки $\gamma = \alpha \beta$):
\[
  \texttt{Half}(L) = \{ \gamma \in \Sigma^\star \mid \gamma \gamma \in L \}.
\]
Разликата е в това $\alpha$, което махаме от $\gamma$, за да получим $\beta$.
Оказва се, че:
\[
  \texttt{Part}(L) = \texttt{Suff}(\texttt{Half}(L))\footnote[2]{Напомням, че $\texttt{Suff}(L) = \{ \gamma \in \Sigma^\star \mid (\exists \omega \in \Sigma^\star)(\omega \gamma \in L) \}$.}.
\]
Нека проверим това:
\begin{align*}
  \beta \in \texttt{Part}(L) \longleftrightarrow & (\exists \alpha \in \Sigma^\star)[\alpha \beta \cdot \alpha \beta \in L] \\ 
                             \longleftrightarrow & (\exists \alpha \in \Sigma^\star)[\alpha \beta \in \texttt{Half}(L)] \\ 
                             \longleftrightarrow & \beta \in \texttt{Suff}(\texttt{Half}(L)).
\end{align*}
Фактът, че $\texttt{Suff}$ запазва регулярност, може да се използва наготово.
Това, че $\texttt{Half}$ запазва регулярност, очакваме да се докаже, но тук ще бъде спестено, понеже е правено на упражнения и го има в лекционните записки.
\end{solution}

\end{document}
